\documentclass[letterpaper]{article}\usepackage[]{graphicx}\usepackage[]{color}
% maxwidth is the original width if it is less than linewidth
% otherwise use linewidth (to make sure the graphics do not exceed the margin)
\makeatletter
\def\maxwidth{ %
  \ifdim\Gin@nat@width>\linewidth
    \linewidth
  \else
    \Gin@nat@width
  \fi
}
\makeatother

\definecolor{fgcolor}{rgb}{0.345, 0.345, 0.345}
\newcommand{\hlnum}[1]{\textcolor[rgb]{0.686,0.059,0.569}{#1}}%
\newcommand{\hlstr}[1]{\textcolor[rgb]{0.192,0.494,0.8}{#1}}%
\newcommand{\hlcom}[1]{\textcolor[rgb]{0.678,0.584,0.686}{\textit{#1}}}%
\newcommand{\hlopt}[1]{\textcolor[rgb]{0,0,0}{#1}}%
\newcommand{\hlstd}[1]{\textcolor[rgb]{0.345,0.345,0.345}{#1}}%
\newcommand{\hlkwa}[1]{\textcolor[rgb]{0.161,0.373,0.58}{\textbf{#1}}}%
\newcommand{\hlkwb}[1]{\textcolor[rgb]{0.69,0.353,0.396}{#1}}%
\newcommand{\hlkwc}[1]{\textcolor[rgb]{0.333,0.667,0.333}{#1}}%
\newcommand{\hlkwd}[1]{\textcolor[rgb]{0.737,0.353,0.396}{\textbf{#1}}}%
\let\hlipl\hlkwb

\usepackage{framed}
\makeatletter
\newenvironment{kframe}{%
 \def\at@end@of@kframe{}%
 \ifinner\ifhmode%
  \def\at@end@of@kframe{\end{minipage}}%
  \begin{minipage}{\columnwidth}%
 \fi\fi%
 \def\FrameCommand##1{\hskip\@totalleftmargin \hskip-\fboxsep
 \colorbox{shadecolor}{##1}\hskip-\fboxsep
     % There is no \\@totalrightmargin, so:
     \hskip-\linewidth \hskip-\@totalleftmargin \hskip\columnwidth}%
 \MakeFramed {\advance\hsize-\width
   \@totalleftmargin\z@ \linewidth\hsize
   \@setminipage}}%
 {\par\unskip\endMakeFramed%
 \at@end@of@kframe}
\makeatother

\definecolor{shadecolor}{rgb}{.97, .97, .97}
\definecolor{messagecolor}{rgb}{0, 0, 0}
\definecolor{warningcolor}{rgb}{1, 0, 1}
\definecolor{errorcolor}{rgb}{1, 0, 0}
\newenvironment{knitrout}{}{} % an empty environment to be redefined in TeX

\usepackage{alltt}\usepackage[]{graphicx}\usepackage[]{color}
\usepackage{framed}
\usepackage{alltt}
\pdfpagewidth 8.5in
\pdfpageheight 11.0in

\usepackage{ucs}
\usepackage[utf8x]{inputenc}
\newcommand{\HRule}{\rule{\linewidth}{0.5mm}}


\usepackage[hmargin=3cm,vmargin=3.5cm]{geometry}
\usepackage{ucs}
\usepackage[utf8x]{inputenc}
\usepackage{amsmath}
\usepackage{amsfonts}
\usepackage{dcolumn}
\usepackage{fancyhdr}
\usepackage{amssymb}
\usepackage{hyperref}
\usepackage[justification=centering]{caption}
\usepackage{enumerate}
\usepackage{graphicx}
\usepackage{cancel}
\usepackage{pbox}
\newcounter{exer}
\numberwithin{exer}{section}
\usepackage{color}
\usepackage[round]{natbib}
\bibliographystyle{plainnat}

\pagestyle{fancyplain}
\usepackage{lastpage}
\definecolor{hellgrau2}{gray}{.6}
\renewcommand{\headrulewidth}{0pt}
\renewcommand{\footrulewidth}{1pt \color{hellgrau2}}
\newcommand{\co}{CO\ensuremath{_2} }
\fancyfoot{}
\fancyhf{}
\rfoot{\textbf{Page \thepage\ of \pageref{LastPage}}}
 \lfoot{\textbf{Econ 102}}
 \cfoot{\textbf{Module 7: Formulas}}

\usepackage{mdframed}
\usepackage{ntheorem}
\theoremstyle{nonumberplain}
\theorembodyfont{\upshape}
\definecolor{bisque}{rgb}{1.0, 0.89, 0.77}
\definecolor{etonblue}{rgb}{0.59, 0.78, 0.64}
\definecolor{blizzardblue}{rgb}{0.67, 0.9, 0.93}
\definecolor{cornsilk}{rgb}{1.0, 0.97, 0.86}
\definecolor{banana}{rgb}{0.98, 0.91, 0.71}
\newmdtheoremenv[leftmargin=20, rightmargin=20,backgroundcolor=bisque,ntheorem]{background}{Optional background Information}
\newmdtheoremenv[leftmargin=20, rightmargin=20,backgroundcolor=etonblue,ntheorem]{important}{Important}
\newmdtheoremenv[leftmargin=20, rightmargin=20,backgroundcolor=cornsilk,ntheorem]{anime}{Suggested Animation}
\newmdtheoremenv[leftmargin=20, rightmargin=20,backgroundcolor=banana,ntheorem]{pause}{Pause \& Practice}
\newmdtheoremenv[leftmargin=20, rightmargin=20,backgroundcolor=bisque,ntheorem]{devel}{Note to the developers}
\usepackage{soul}
\newcommand{\mathcolorbox}[1]{\colorbox{blizzardblue}{$\displaystyle #1$}}
\IfFileExists{upquote.sty}{\usepackage{upquote}}{}
\IfFileExists{upquote.sty}{\usepackage{upquote}}{}
\begin{document}
\noindent
 \textsc{\Large ECON 102: Inequality Formulas}\\[.1cm]
 \begin{center}
   \line(1,0){450}\\[0.2cm]
   \end{center}

In this document, we present the most important formulas used in
Module 7. 

\section*{Income Share}

To illustrate, suppose the population is composed of 20 individuals
with the following incomes:
\begin{center}
% latex table generated in R 4.0.2 by xtable 1.8-4 package
% Sun Jul 12 15:32:00 2020
\begin{tabular}{rrrrrrrrrr}
  \hline
  \hline
 39 & 104 &  76 & 193 & 114 & 134 & 183 &  61 &  75 &  99 \\ 
  143 & 103 &  32 & 121 & 111 &  54 & 130 & 117 & 152 &  57 \\ 
   \hline
\end{tabular}

\end{center}
The first step is to sort the income:
\begin{center}
% latex table generated in R 4.0.2 by xtable 1.8-4 package
% Sun Jul 12 15:32:00 2020
\begin{tabular}{rrrrrrrrrr}
  \hline
  \hline
 32 &  39 &  54 &  57 &  61 &  75 &  76 &  99 & 103 & 104 \\ 
  111 & 114 & 117 & 121 & 130 & 134 & 143 & 152 & 183 & 193 \\ 
   \hline
\end{tabular}

\end{center}

\subsection*{Income Share by Quintile}

We group the population in five equally populated groups from the
poorest to the richest and we compute the total income in each group:

\begin{center}
% latex table generated in R 4.0.2 by xtable 1.8-4 package
% Sun Jul 12 15:32:00 2020
\begin{tabular}{lcccc||c}
  \hline
  &&&&&Total\\ \hline
Quintile 1 &  32 &  39 &  54 &  57 & 182 \\ 
  Quintile 2 &  61 &  75 &  76 &  99 & 311 \\ 
  Quintile 3 & 103 & 104 & 111 & 114 & 432 \\ 
  Quintile 4 & 117 & 121 & 130 & 134 & 502 \\ 
  Quintile 5 & 143 & 152 & 183 & 193 & 671 \\ 
   \hline
\end{tabular}

\end{center}
Then, we divide the total of each group by the total income. For
example, the income share in the first quintile is 
\[
\frac{182}{182+311+432+502+671}\times 100 = 
8.6749\%
\]
The income shares by quintile in percentage are therefore:
\begin{center}
% latex table generated in R 4.0.2 by xtable 1.8-4 package
% Sun Jul 12 15:32:00 2020
\begin{tabular}{rr}
  \hline
  \hline
Quintile 1 & 8.6749 \\ 
  Quintile 2 & 14.8236 \\ 
  Quintile 3 & 20.5910 \\ 
  Quintile 4 & 23.9276 \\ 
  Quintile 5 & 31.9828 \\ 
   \hline
\end{tabular}

\end{center}

\subsubsection*{Cumulative Income Share by Quintile}

The cumulative sum of a list of numbers $\{X_1, ..., X_n\}$ is the following list:
\begin{eqnarray*}
&&X_1\\
&&X_1+X_2\\
&&X_1+X_2+X_3\\
&&X_1+X_2+X_3+X_4\\
&& \vdots\\
&& X_1+X_2+X_3+X_4+\cdots+X_n
\end{eqnarray*}

Apply this definition to the income shares and we obtain the
cumulative income shares. The second column is the cumulative
population share (20\% of population are in each quintile).
\begin{center}
% latex table generated in R 4.0.2 by xtable 1.8-4 package
% Sun Jul 12 15:32:00 2020
\begin{tabular}{rrr}
  \hline
 & Cum. income share & Population Share \\ 
  \hline
Quintile 1 & 8.7 & 20.0 \\ 
  Quintile 2 & 23.5 & 40.0 \\ 
  Quintile 3 & 44.1 & 60.0 \\ 
  Quintile 4 & 68.0 & 80.0 \\ 
  Quintile 5 & 100.0 & 100.0 \\ 
   \hline
\end{tabular}

\end{center}

\subsubsection*{Gini Coefficient From Quintiles}

Let $h_i$ be the cumulative income share for quintile i not expressed
in percentage, then the Gini coefficient is:
\[
G = 0.8 - 0.4\times(h_1+h_2+h_3+h_4)
\]
In this case, it is equal to 0.2228.

\subsection*{Income Share by Decile}

We group the population in 10 equally populated groups from the poorest
to the richest and we compute the total income in each group:

\begin{center}
% latex table generated in R 4.0.2 by xtable 1.8-4 package
% Sun Jul 12 15:32:00 2020
\begin{tabular}{lcc||c}
  \hline
  &&&Total\\ \hline
Decile 1 &  32 &  39 &  71 \\ 
  Decile 2 &  54 &  57 & 111 \\ 
  Decile 3 &  61 &  75 & 136 \\ 
  Decile 4 &  76 &  99 & 175 \\ 
  Decile 5 & 103 & 104 & 207 \\ 
  Decile 6 & 111 & 114 & 225 \\ 
  Decile 7 & 117 & 121 & 238 \\ 
  Decile 8 & 130 & 134 & 264 \\ 
  Decile 9 & 143 & 152 & 295 \\ 
  Decile 10 & 183 & 193 & 376 \\ 
   \hline
\end{tabular}

\end{center}

Then, we divide the total of each group by the total income as for
the quintile case. The income shares by decile in percentage are
therefore:
\begin{center}
% latex table generated in R 4.0.2 by xtable 1.8-4 package
% Sun Jul 12 15:32:00 2020
\begin{tabular}{lc}
  \hline
  \hline
Decile 1 & 3.3842 \\ 
  Decile 2 & 5.2908 \\ 
  Decile 3 & 6.4824 \\ 
  Decile 4 & 8.3413 \\ 
  Decile 5 & 9.8665 \\ 
  Decile 6 & 10.7245 \\ 
  Decile 7 & 11.3441 \\ 
  Decile 8 & 12.5834 \\ 
  Decile 9 & 14.0610 \\ 
  Decile 10 & 17.9218 \\ 
   \hline
\end{tabular}

\end{center}

\subsubsection*{Cumulative Income Share by Decile}

We apply the same formula presented earlier for the quintiles:

\begin{center}
% latex table generated in R 4.0.2 by xtable 1.8-4 package
% Sun Jul 12 15:32:00 2020
\begin{tabular}{lcc}
  \hline
 & Cum. income share & Population Share \\ 
  \hline
Decile 1 & 3.4 & 10.0 \\ 
  Decile 2 & 8.7 & 20.0 \\ 
  Decile 3 & 15.2 & 30.0 \\ 
  Decile 4 & 23.5 & 40.0 \\ 
  Decile 5 & 33.4 & 50.0 \\ 
  Decile 6 & 44.1 & 60.0 \\ 
  Decile 7 & 55.4 & 70.0 \\ 
  Decile 8 & 68.0 & 80.0 \\ 
  Decile 9 & 82.1 & 90.0 \\ 
  Decile 10 & 100.0 & 100.0 \\ 
   \hline
\end{tabular}

\end{center}

\subsubsection*{Gini Coefficient From Quintiles}

Let $h_i$ be the cumulative income share for decile i, then the Gini
coefficient is:
\[
G = 0.9 - 0.2\times(h_1+h_2+h_3+h_4+h_5+h_6+h_7+h_8+h_9)
\]
In this case, it is equal to 0.2324.

\section*{The Median}

The median is the number that separates a sample in two equal groups
after sorting the observations.

\subsubsection*{Odd Number of Observations}

If the series is $\{X_1, X_2, X_3, X_4, X_5\}$ with
$X_1<X_2<X_3<X_4<X_5$, the median is equal to $X_3$.

\subsubsection*{Even Number of Observations}

If the series is $\{X_1, X_2, X_3, X_4, X_5, X_6\}$ with
$X_1<X_2<X_3<X_4<X_5<X_6$, the median (one version of it) is equal to
$(X_3+X_4)/2$.



\end{document}
